% from ega4-8-fr.tex
\oldpage[IV-3]{17}
alors, pour que $Z$ soit ouverte (resp. fermée, resp. localement fermée) dans
$S$, il faut et il suffit qu'il existe $\lambda\geq\alpha$ tel que
$Z_\lambda$ le soit dans $S_\lambda$.
\end{corollary}

Couvrons $S_\alpha$ par un nombre fini d'ouverts affines
$U_\alpha^{(i)}$~; alors les $U^{(i)}=u_\alpha^{-1}(U_\alpha^{(i)})$ forment
un recouvrement ouvert affine de $S$, et pour que $Z$ soit ouvert (resp.
fermé, resp. localement fermé) dans $S$, il faut et il suffit que chacun des
$Z\cap U^{(i)}$ le soit dans $U^{(i)}$. Comme $L$ est filtrant et que chacun
des $Z\cap U^{(i)}$ est constructible dans $U^{(i)}$
$(0_{\mathrm{III}},9.1.8)$, on peut se borner à démontrer le corollaire
lorsque $S_\alpha$ est affine, donc quasi-compact et quasi-séparé~; mais alors
il résulte de (8.3.11).

\begin{proposition}[8.3.13]
\label{IV.8.3.13-fr}
Supposons que les morphismes $u_{\lambda\mu}:S_\mu\to S_\lambda$ soient
plats pour $\lambda\leq\mu$, et qu'il existe $\alpha$ tel que $S_\alpha$ soit
quasi-compact. Pour tout $\lambda$, soient $Z'_\lambda$, $Z''_\lambda$ deux
parties pro-constructibles de $S_\lambda$, telles que
$Z'_\mu=u_{\lambda\mu}^{-1}(Z'_\lambda)$ et
$Z''_\mu=u_{\lambda\mu}^{-1}(Z''_\lambda)$ pour $\mu\geq\lambda$~; on
suppose en outre que $\overline{Z'_\alpha}$ est constructible dans
$S_\alpha$. Soient $Z'=u_\lambda^{-1}(Z'_\lambda)$,
$Z''=u_\lambda^{-1}(Z''_\lambda)$~; pour que
$Z''\subset\overline{Z'}$, il faut et il suffit qu'il existe
$\lambda\geq\alpha$ tel que
$Z''_\lambda\subset\overline{Z'_\lambda}$.
\end{proposition}

En effet, on sait que $u_\lambda:S\to S_\lambda$ est aussi un morphisme plat
pour tout $\lambda$ (8.3.8)~; comme $Z'_\lambda$ est pro-constructible,
l'adhérence de $Z'_\mu$ pour $\mu\geq\lambda$ (resp. de $Z'$) dans $S_\mu$
(resp. $S$) est égale à
$u_{\lambda\mu}^{-1}(\overline{Z'_\lambda})$ (resp.
$u_\lambda^{-1}(\overline{Z'_\lambda})$) (2.3.10). Comme les
$u_{\lambda\mu}^{-1}(\overline{Z'_\lambda})$ et
$u_\lambda^{-1}(\overline{Z'_\lambda})$ sont constructibles (1.8.2), la
conclusion résulte de (8.3.2).

\subsection{Critères d'irréductibilité et de connexion pour les limites
projectives de préschémas}
\label{subsection:IV.8.4-fr}

\begin{proposition}[8.4.1]
\label{IV.8.4.1-fr}
Supposons qu'il existe un indice $\alpha$ tel que $S_\alpha$ soit
quasi-compact.
\begin{enumerate}
\item[(i)] Si $S$ n'est pas irréductible et si de plus l'espace sous-jacent à
$S$ est noethérien et $S_\alpha$ quasi-séparé, il existe
$\lambda\geq\alpha$ tel que, pour $\mu\geq\lambda$, $S_\mu$ ne soit pas
irréductible.
\item[(ii)] Si $S$ n'est pas connexe, il existe $\lambda\geq\alpha$ tel que,
pour $\mu\geq\lambda$, $S_\mu$ ne soit pas connexe.
\end{enumerate}
\end{proposition}

Supposons que $S$ soit réunion de deux ensembles fermés $S'$, $S''$ distincts
de $S$ (resp. fermés disjoints et non vides). Dans le cas (i), $S'$ et $S''$
sont constructibles puisque l'espace $S$ est noethérien. En vertu de
(8.3.11), il existe donc un $\lambda\geq\alpha$ et deux ensembles fermés
constructibles $S'_\lambda$, $S''_\lambda$ de $S_\lambda$ tels que
$S'=u_\lambda^{-1}(S'_\lambda)$,
$S''=u_\lambda^{-1}(S''_\lambda)$~; comme $S=S'\cup S''$, il résulte aussi
de (8.3.11) que l'on peut supposer que
$S_\lambda=S'_\lambda\cup S''_\lambda$~; comme $S'_\lambda$ et
$S''_\lambda$ sont distincts de $S_\lambda$, cela prouve que $S_\lambda$
n'est pas irréductible.

Dans le cas (ii), $S'$ et $S''$ sont ouverts quasi-compacts, donc, en vertu
de (8.2.11), il existe un $\lambda\geq\alpha$ et deux ensembles ouverts
quasi-compacts $S'_\lambda$, $S''_\lambda$ de $S_\lambda$ tels que
$S'=u_\lambda^{-1}(S'_\lambda)$,
$S''=u_\lambda^{-1}(S''_\lambda)$. D'ailleurs, comme $S'$ et $S''$ sont
ouverts et fermés dans $S$, ils sont à la fois pro-constructibles et
ind-constructibles (1.9.6), donc constructibles (1.9.11), et il résulte donc
de (8.3.5) que l'on peut supposer $\lambda$ pris tel que
$S_\lambda=S'_\lambda\cup S''_\lambda$ et
$S'_\lambda\cap S''_\lambda=\emptyset$, ce qui montre que $S_\lambda$
n'est pas connexe.

\begin{proposition}[8.4.2]
\label{IV.8.4.2-fr}
Supposons que l'espace sous-jacent à $S$ soit noethérien et que l'une des deux
conditions suivantes soit vérifiée~:
\begin{enumerate}
\item[\emph{a)}] Pour $\lambda\leq\mu$,
$u_{\lambda\mu}:S_\mu\to S_\lambda$ est dominant, et il existe $\alpha$ tel
que $S_\alpha$ soit quasi-compact.
